\abstracttext{
Mars presents a highly anomalous magnetic environment characterized by the extinction of its ancient core dynamo and the survival of intense, localized crustal magnetic remanence predominantly in its southern ancient cratered highlands. As international space agencies and commercial ventures plan crewed missions and long-term surface habitats on Mars, understanding this complex magnetic landscape is vital for designing bioregenerative life support systems (BLSS), agricultural growth chambers, and astronaut health protocols. In this study, we quantitatively map the Martian crustal magnetic field across 13 contemporary, historical, and candidate human landing sites using calibrated spherical harmonic models derived from Mars Global Surveyor (MGS) and MAVEN orbital magnetometer observations, integrated with InSight surface ground truth. We demonstrate that prospective human habitat sites situated in the ice-rich northern plains (e.g., Arcadia Planitia and Deuteronilus Mensae) and rover landing sites (e.g., Jezero Crater and Oxia Planum) reside in near-null hypomagnetic field (HMF) environments ($|\mathbf{B}| < \SI{20}{\nano\tesla}$ at surface level, $>2,500\times$ weaker than Earth's $\sim\SI{50}{\micro\tesla}$ geomagnetic field). Conversely, southern highland anomaly belts (e.g., Terra Sirenum and Terra Cimmeria) exhibit ground-level field intensities exceeding \SI{1500}{\nano\tesla}, generating localized mini-magnetospheres with distinct plasma deflection capabilities. We synthesize the experimental magnetobiology literature across plant growth, root directional auxin transport, cryptochrome signaling, microbial community stability, and mammalian bone demineralization, formulating a biological risk matrix and operational siting framework for Martian surface exploration.
}

\keywordstext{Mars crustal magnetic fields, Hypomagnetic field (HMF), InSight magnetometer, Space biology, Bioregenerative life support, Plant morphogenesis, Human Mars exploration}
